\documentclass[11pt,reqno]{amsart}

\usepackage{amsmath,amssymb,amsthm}
\usepackage[T1]{fontenc}
\usepackage[expansion=false]{microtype}
\usepackage[margin=1in]{geometry}
\usepackage{booktabs}
\usepackage{hyperref}
\hypersetup{colorlinks=true,linkcolor=blue,citecolor=blue,urlcolor=blue}

\theoremstyle{plain}
\newtheorem{theorem}{Theorem}[section]
\newtheorem{proposition}[theorem]{Proposition}
\newtheorem{lemma}[theorem]{Lemma}
\newtheorem{corollary}[theorem]{Corollary}

\theoremstyle{definition}
\newtheorem{definition}[theorem]{Definition}
\newtheorem{observation}[theorem]{Observation}
\newtheorem{conjecture}[theorem]{Conjecture}

\theoremstyle{remark}
\newtheorem{remark}[theorem]{Remark}

\newcommand{\Q}{\mathbb{Q}}
\newcommand{\Z}{\mathbb{Z}}
\newcommand{\F}{\mathbb{F}}
\newcommand{\PCP}{E_{\mathrm{PCP}}}
\DeclareMathOperator{\rank}{rank}
\DeclareMathOperator{\Reg}{Reg}
\DeclareMathOperator{\disc}{disc}
\DeclareMathOperator{\NS}{NS}

\sloppy

\begin{document}

\title[Mordell--Weil ranks and a survey of the perfect-cuboid elliptic family]{Mordell--Weil ranks and an experimental survey of the
perfect-cuboid elliptic family}

\author{Lightman Chang}
\address{Independent Researcher}
\email{lightman.chang@gmail.com}

\subjclass[2020]{Primary 11G05; Secondary 11D09, 11G07, 11Y50, 14J27}
\keywords{perfect cuboid, Euler brick, elliptic surface, Mordell--Weil rank,
rank distribution, computational number theory}

\date{\today}

\begin{abstract}
A perfect cuboid is a rectangular box whose three edges, three face diagonals,
and space diagonal are all integers; whether one exists is a question of Euler
(1769) that remains open. To each Pythagorean parameter pair $(m,n)$ one
attaches the elliptic fiber $\PCP(q)\colon y^2=x(x+1)(x+q^2)$, where
$q=(m^2-n^2)/(2mn)$, and the existence of a perfect cuboid forces a rational
point on a finite list of such fibers. We present three computational
contributions, all carried out in PARI/GP~2.15.4 with exact (square-test and
$2$-descent) arithmetic. First, we record a multi-framework, machine-checked
confirmation that no perfect cuboid arises among the $36$ primitive Euler
bricks with all edges at most $30{,}000$. Second, we tabulate the certified
Mordell--Weil rank distribution across $303$ primitive Pythagorean fibers
($2\le m\le 38$): ranks $0,1,2,3$ occur with relative frequencies
$38.9\%,45.2\%,14.9\%,1.0\%$, with maximum observed rank $3$ and no fiber left
rank-ambiguous by $2$-descent. Third, and as the single proved theorem, we
certify that the fiber attached to $(m,n)=(22,17)$, namely $q=195/748$, has
Mordell--Weil rank exactly $3$ over $\Q$, by exhibiting three explicit rational
points whose canonical-height regulator equals $18.8337\ldots\neq 0$ together
with the matching upper bound from $2$-descent; this complements the generic
(geometric) rank statements of Naskr\k{e}cki by an explicit high-rank
$\Q$-fiber. We close with a short catalogue of the named curves arising in the
problem's reformulations, distinguishing throughout between theorem,
observation, and conjecture.
\end{abstract}

\maketitle

\section{Introduction}\label{sec:intro}

The \emph{Perfect Cuboid Problem} (PCP), raised by Euler in 1769, asks whether
there exist positive integers $a,b,c$ such that the three face diagonals
$\sqrt{a^2+b^2}$, $\sqrt{b^2+c^2}$, $\sqrt{a^2+c^2}$ and the space diagonal
$\sqrt{a^2+b^2+c^2}$ are all integers. A triple $(a,b,c)$ satisfying only the
three face conditions is an \emph{Euler brick}; these exist in abundance (the
smallest is $(44,117,240)$). No perfect cuboid is known, and none is ruled out;
the problem is recorded as Problem~D18 in Guy~\cite{Guy}. We refer to
van~Luijk~\cite{vanLuijk}, Sharipov~\cite{Sharipov}, Naskr\k{e}cki~\cite{Nas},
Yoshida~\cite{Yoshida}, and Peschmann~\cite{Pesch1,Pesch2} for the modern
arithmetic-geometry viewpoint, which we adopt.

\subsection*{The per-fiber elliptic family}
A standard route attaches to the problem the elliptic surface
\[
  \PCP(q)\colon\qquad y^2 = x(x+1)(x+q^2),
\]
parametrized by a rational $q$ that runs over leg-ratios of Pythagorean
triples. Concretely, for coprime opposite-parity $m>n\ge 1$ set
\begin{equation}\label{eq:param}
  a = m^2-n^2,\qquad b = 2mn,\qquad q = \frac{a}{b}.
\end{equation}
The substitution $x\mapsto x/b^2$, $y\mapsto y/b^3$ identifies $\PCP(q)$ with
the integral model
\begin{equation}\label{eq:Emn}
  E(m,n)\colon\quad y^2 = x(x+b^2)(x+a^2) = x^3 + (a^2+b^2)\,x^2 + (ab)^2\,x,
\end{equation}
which has good arithmetic for computation. Over the function field $\Q(q)$ the
generic curve has Mordell--Weil rank $1$ (and $2$ on an explicit subfamily), as
proved by Naskr\k{e}cki~\cite{Nas} for the closely related family
$y^2=x(x-a^2)(x-b^2)$; this is a statement about the \emph{geometric} or
\emph{generic} rank. The present paper is concerned instead with the
\emph{arithmetic} ranks of \emph{individual} $\Q$-fibers and with their
empirical distribution.

\subsection*{What this paper does}
This is a paper of experimental mathematics. Its content is computational, and
we are careful to label each statement by its logical status:
\begin{itemize}
  \item a \emph{Theorem} is a statement we prove (here, exactly one);
  \item an \emph{Observation} is a verified numerical fact about a finite,
        explicitly searched range;
  \item a \emph{Conjecture} is an extrapolation beyond the searched range that
        we do not prove.
\end{itemize}
The three components are: a machine-checked no-cuboid record up to edge
$30{,}000$ (Section~\ref{sec:bricks}); the Mordell--Weil rank distribution over
$303$ Pythagorean fibers (Section~\ref{sec:dist}); and an explicit
$\Q$-rank-$3$ fiber, certified as Theorem~\ref{thm:main}
(Section~\ref{sec:rank3}). A short catalogue of named curves
(Section~\ref{sec:catalogue}) records the reformulations used. All scripts and
their outputs are listed in Section~\ref{sec:repro}.

\subsection*{Relation to prior work}
None of the three components claims to make progress toward deciding PCP. The
no-cuboid record extends the explicit verifications surveyed in
\cite{vanLuijk,Sharipov} and is complementary to the conditional and per-fiber
closures of Peschmann~\cite{Pesch1,Pesch2}, who treats $1{,}072$ fibers by a
torsion-intersection method requiring a rank-zero elliptic quotient. The rank
distribution is, to our knowledge, the first tabulation with every fiber
\emph{certified} (no rank gaps) by $2$-descent over a stated range. The
explicit $\Q$-rank-$3$ fiber is the one genuinely new theorem: Naskr\k{e}cki's
optimal lower bounds are generic/geometric (rank $1$ generically, $2$ on a
subfamily), and an individual $\Q$-fiber of rank as high as $3$ is exhibited
here with full certification.

\section{No perfect cuboid among small Euler bricks}\label{sec:bricks}

\subsection*{Algorithm}
We enumerate Euler bricks by a face sieve and test the space-diagonal condition
on each. All tests use PARI's exact integer \texttt{issquare}; no floating-point
comparison enters, so the search is rigorous on its range.

\begin{verbatim}
Input : EDGEMAX.
Output: all primitive Euler bricks with edges <= EDGEMAX; any perfect cuboid.
for a = 1 .. EDGEMAX:
  for b = a .. EDGEMAX:
    if issquare(a^2 + b^2):                     # (a,b) a Pythagorean leg-pair
      for c = b .. EDGEMAX:
        if issquare(b^2 + c^2) and issquare(a^2 + c^2):
          if gcd(a, b, c) == 1:                 # primitive Euler brick
            record (a, b, c)
            if issquare(a^2 + b^2 + c^2):        # space diagonal integral
              report PERFECT CUBOID (a, b, c)
\end{verbatim}
The loop ordering $a\le b\le c$ enumerates each unordered brick once.

\subsection*{Search range and result}
We ran the algorithm with $\mathrm{EDGEMAX}=30{,}000$ (PARI/GP~2.15.4, single
core, $\approx 3$ minutes wall time). The smaller cutoffs $1000$ and $5000$
were run as cross-checks and reproduce the brick counts $5$ and $11$ recorded
in the literature.

\begin{observation}\label{obs:bricks}
There are exactly $36$ primitive Euler bricks with all edges at most $30{,}000$,
and none of them is a perfect cuboid: for every one of the $36$ triples
$(a,b,c)$, the integer $a^2+b^2+c^2$ is not a perfect square.
\end{observation}

The complete list of the $36$ bricks $(a\le b\le c)$ found is
\[
\begin{array}{llll}
(44,117,240) & (85,132,720) & (140,480,693) & (160,231,792) \\
(187,1020,1584) & (195,748,6336) & (240,252,275) & (429,880,2340) \\
(495,4888,8160) & (528,5796,6325) & (780,2475,2992) & (828,2035,3120) \\
(832,855,2640) & (935,17472,25704) & (1008,1100,1155) & (1008,1100,12075) \\
(1080,1881,14560) & (1155,6300,6688) & (1560,2295,5984) & (1575,1672,9120) \\
(1755,4576,6732) & (2925,3536,11220) & (2964,9152,9405) & (4368,4901,13860) \\
(4599,18368,23760) & (4900,17157,23760) & (5643,14160,21476) & (6072,16929,18560) \\
(6435,24080,24684) & (7579,8820,17472) & (7800,23751,29920) & (7840,9828,10725) \\
(7920,15232,26649) & (8789,10560,17748) & (10296,11753,16800) & (14112,15400,19305).
\end{array}
\]

\begin{remark}
Observation~\ref{obs:bricks} is a finite statement about a bounded range, not a
nonexistence theorem. We state it as an observation and not a theorem precisely
to keep that distinction explicit.
\end{remark}

\section{Mordell--Weil rank distribution across Pythagorean fibers}\label{sec:dist}

\subsection*{Sample and method}
We take every primitive Pythagorean parameter pair $(m,n)$ with
\[
  2\le m\le 38,\qquad 1\le n<m,\qquad \gcd(m,n)=1,\qquad m\not\equiv n \pmod 2,
\]
forming the integral model $E(m,n)$ of \eqref{eq:Emn} and its minimal model via
PARI's \texttt{ellminimalmodel}. For each fiber we call \texttt{ellrank} at
effort level $1$, obtaining an interval $[\ell,h]$ where $\ell$ is a certified
lower bound (witnessed by explicit independent points) and $h$ is the upper
bound from $2$-descent. We record a \emph{certified} rank only when
$\ell=h$, and tally an \emph{uncertified} bucket otherwise; we never guess a
rank when $2$-descent leaves a gap.

\begin{observation}\label{obs:dist}
Among the $303$ primitive Pythagorean fibers in the range $2\le m\le 38$, every
fiber has $\ell=h$ (no uncertified fibers), and the certified rank distribution
is
\begin{center}
\begin{tabular}{rrr}
\toprule
rank & count & frequency \\
\midrule
$0$ & $118$ & $38.94\%$ \\
$1$ & $137$ & $45.21\%$ \\
$2$ & $45$ & $14.85\%$ \\
$3$ & $3$ & $0.99\%$ \\
\midrule
total & $303$ & $100\%$ \\
\bottomrule
\end{tabular}
\end{center}
The maximum rank observed is $3$, attained by exactly three fibers, namely
those with $(m,n)\in\{(22,17),(35,22),(37,26)\}$.
\end{observation}

The combined $\rank\ge 2$ frequency is $15.8\%$, which is larger than a naive
$50/50$ parity heuristic would predict and which increases with $m$ over the
range. This behaviour is consistent with the rank distribution expected for a
family whose conductor grows with the parameters, but we make no quantitative
claim and isolate the qualitative trend as a conjecture only.

\begin{conjecture}\label{conj:rank}
For the family $\PCP(q)$ with $q$ as in \eqref{eq:param}, rank $3$ fibers occur
with positive frequency as $m\to\infty$, and the maximum rank over
$2\le m\le M$ is unbounded in $M$.
\end{conjecture}

We do not have evidence beyond $m=38$ for the unboundedness assertion, and
flag it as speculative; the project literature reports rank-$4$ fibers at larger
$m$ obtained by a targeted (non-exhaustive) high-$\omega$ search, but those lie
outside the exhaustive range surveyed here and we do not import them.

\section{An explicit fiber of rank three over \texorpdfstring{$\Q$}{Q}}\label{sec:rank3}

The three rank-$3$ fibers of Observation~\ref{obs:dist} are the smallest
high-rank specializations in our range. We single out the first,
$(m,n)=(22,17)$, and certify its rank rigorously. Here
$a=195$, $b=748$, $q=195/748$, and the minimal model of $E(22,17)$ is
\begin{equation}\label{eq:E2217}
  y^2 + xy = x^3 - 6{,}108{,}655{,}980\,x + 180{,}712{,}439{,}349{,}327,
\end{equation}
with conductor $N = 19{,}015{,}731{,}735 = 3\cdot5\cdot7\cdot11\cdot13\cdot17\cdot23\cdot41\cdot79$ and torsion subgroup $\Z/4\times\Z/2$.

\begin{theorem}\label{thm:main}
The elliptic curve $\PCP(q)$ with $q=195/748$ has Mordell--Weil rank exactly
$3$ over $\Q$. The three rational points
\[
  Q_1 = \Big(-\tfrac{15}{176},\ \tfrac{2415}{65824}\Big),\quad
  Q_2 = \Big(\tfrac{117}{44},\ \tfrac{169533}{32912}\Big),\quad
  Q_3 = \Big(\tfrac{12675}{44},\ \tfrac{161213325}{32912}\Big)
\]
on the curve
\[
  \PCP(195/748)\colon\quad y^2=x(x+1)\bigl(x+(195/748)^2\bigr)
\]
are independent in
$\PCP(195/748)(\Q)/\PCP(195/748)(\Q)_{\mathrm{tors}}$, and they span a
finite-index subgroup of the free part.
\end{theorem}

\begin{proof}
We work on the integral model $E(22,17)$ of \eqref{eq:Emn}; the substitution
$x\mapsto x/b^2$, $y\mapsto y/b^3$ with $b=748$ carries its points to those of
$\PCP(q)$ and is a $\Q$-isomorphism, so it suffices to certify rank $3$ for
$E(22,17)$ and to transport three generators.

\emph{Upper bound.} PARI's \texttt{ellrank} performs a complete $2$-descent on
the integral model $E(22,17)$ of \eqref{eq:Emn}, using the full rational
$2$-torsion (the $2$-Selmer rank bound is isomorphism-invariant, so the choice
between $E(22,17)$ and its reduced minimal model \eqref{eq:E2217} is immaterial).
It returns the interval $[\ell,h]=[3,3]$; the value $h=3$ is an unconditional
upper bound on the rank coming from the size of the $2$-Selmer group (no analytic
or $\mathrm{Sha}$-finiteness hypothesis is used for the upper bound).

\emph{Lower bound.} The same call returns three independent rational points on
$E(22,17)$; transporting them by the $\Q$-isomorphism $x\mapsto x/b^2$,
$y\mapsto y/b^3$ ($b=748$) gives $Q_1,Q_2,Q_3$ on $\PCP(q)$ as displayed, each
verified to satisfy the curve equation by direct substitution. The
canonical-height pairing matrix of $(Q_1,Q_2,Q_3)$, computed by
\texttt{ellheightmatrix}, is
\[
H = \begin{pmatrix}
2.73288 & -0.50982 & 0.75786 \\
-0.50982 & 2.76342 & 0.84518 \\
0.75786 & 0.84518 & 3.15765
\end{pmatrix},
\]
with regulator $\Reg = \det H = 18.83372\ldots$. Since $\Reg\neq 0$, the three
points are linearly independent in the Mordell--Weil group modulo torsion;
hence the rank is at least $3$.

Combining the two bounds gives rank exactly $3$, and the three independent
points then span a finite-index subgroup of the free part (the index is the
saturation, which we do not compute; it is irrelevant to the rank and to
independence).
\end{proof}

\begin{remark}
The proof is a genuine theorem and not merely an \texttt{ellrank} lower bound:
the certification rests on (i) the unconditional $2$-Selmer upper bound $h=3$
and (ii) the explicit nonzero regulator, both of which are exact-up-to-rounding
PARI outputs that any reader can reproduce or recompute by hand. No conjecture
(BSD, parity, ABC) enters. The displayed regulator is a real-analytic quantity
computed to $30$ significant digits; only its non-vanishing, comfortably away
from zero, is used.
\end{remark}

\begin{remark}[Relation to Naskr\k{e}cki~\cite{Nas}]
Naskr\k{e}cki establishes that the \emph{generic} fiber of the Pythagorean
family has Mordell--Weil rank $1$ over $\Q(q)$ ($2$ on an explicit subfamily),
and that these geometric bounds are optimal. Theorem~\ref{thm:main} is a
statement of a different kind: it pins down an \emph{individual} $\Q$-fiber
whose arithmetic rank, $3$, exceeds the generic rank by the contribution of two
extra non-generic sections. Such specializations are expected to be sparse, in
keeping with Observation~\ref{obs:dist}.
\end{remark}

\section{Named curves in the PCP reformulations: a catalogue}\label{sec:catalogue}

For the reader's convenience we record, with explicit identifications, the
elliptic curves that recur in the problem's reformulations. The general shape of
these reductions --- attaching elliptic quotients to the quartic conditions of
the cuboid --- is folklore, going back through
\cite{vanLuijk,Sharipov,Nas,Pesch1}; we claim originality only for the specific
Cremona-label identifications, which we verified with PARI's
\texttt{ellidentify} and \texttt{ellrank}.

\begin{observation}\label{obs:catalogue}
The following identifications hold.
\begin{enumerate}
  \item The universal fiber $\PCP(q)\colon y^2=x(x+1)(x+q^2)$ has, in the
        variable $t=q^2$, discriminant $\Delta = 16\,t^2(t-1)^2 = 16\,q^4(q^2-1)^2$
        and $c_4 = 16(t^2-t+1)=16(q^4-q^2+1)$, with rational $2$-torsion of
        order $4$ refining to $\Z/4\times\Z/2$; the corresponding rational
        elliptic surface has singular fibers of types
        $I_4, I_2, I_2, I_4$ over $q^2\in\{0,1,\infty\}$ and one further place,
        and Euler number $12$, so its generic Mordell--Weil rank is $0$ over
        $\overline{\Q}(q^2)$.
  \item The ``cleanest formulation'' curve $y^2=x^3+x^2-x+15$, which arises in
        the Saunderson sub-family reduction, is Cremona \texttt{160a2} and has
        rank $1$, with a generator $(-1,4)$.
  \item The auxiliary cover $y^2=x^3-36x^2+320x$ is Cremona \texttt{80a1} and
        has rank $0$.
\end{enumerate}
\end{observation}

The juxtaposition of (ii) and (iii) records the structural feature exploited in
the project's Saunderson-slice closure: a rank-$1$ curve governs the squareness
condition while a rank-$0$ cover bounds the relevant rational points. We state
these as identifications only and make no closure claim here.

\section{Reproducibility}\label{sec:repro}

All computations were performed in PARI/GP version 2.15.4 (GMP kernel) on a
single x86-64 core, with \texttt{parisize} set between $2$ and $6$ GB. The
scripts and their captured outputs are:
\begin{itemize}
  \item \texttt{01\_fiber\_22\_17\_rank.gp} / \texttt{.out} ---
        Theorem~\ref{thm:main}: minimal model, conductor, torsion,
        \texttt{ellrank} interval, height matrix, regulator.
  \item \texttt{02\_rank\_distribution.gp} / \texttt{.out} ---
        Observation~\ref{obs:dist}: the certified rank histogram over the
        $303$ fibers and the list of rank-$3$ fibers.
  \item \texttt{03\_euler\_brick\_search.gp} / \texttt{.out} ---
        Observation~\ref{obs:bricks}: the face-sieve enumeration of the $36$
        primitive Euler bricks with edges $\le 30{,}000$ and the zero
        perfect-cuboid count.
  \item \texttt{04\_named\_curves.gp} / \texttt{.out} ---
        Observation~\ref{obs:catalogue}: the universal geometry and the
        Cremona-label identifications, plus the $q$-fiber coordinates of the
        three generators of Theorem~\ref{thm:main}.
\end{itemize}
The \texttt{ellrank} results are stable across effort levels $0$ and $1$; the
regulator is reported to $30$ significant digits and its non-vanishing is
robust. The brick search uses exact integer square tests throughout.

\section*{Acknowledgements}
The author thanks the maintainers of PARI/GP, the only software used in this
work.

\begin{thebibliography}{9}

\bibitem{Guy}
R.~K.~Guy,
\emph{Unsolved Problems in Number Theory}, 3rd ed.,
Problem Books in Mathematics, Springer, New York, 2004 (Problem~D18).

\bibitem{Nas}
B.~Naskr\k{e}cki,
\emph{Mordell--Weil ranks of families of elliptic curves associated to
Pythagorean triples},
Acta Arith. \textbf{160} (2013), no.~2, 159--183;
arXiv:\allowbreak 1210.6933.

\bibitem{Pesch1}
R.~Peschmann,
\emph{Quartic reductions and elliptic obstructions for perfect Euler bricks},
preprint, 2026; arXiv:\allowbreak 2604.09328.

\bibitem{Pesch2}
R.~Peschmann,
\emph{A torsion-intersection proof of perfect-cuboid nonexistence on $1{,}072$
explicit master-tuple fibers},
preprint, 2026; arXiv:\allowbreak 2604.28072.

\bibitem{Sharipov}
R.~A.~Sharipov,
\emph{On two elliptic curves associated with perfect cuboids},
preprint, 2012; arXiv:\allowbreak 1208.1227.

\bibitem{vanLuijk}
R.~van~Luijk,
\emph{On perfect cuboids},
undergraduate thesis / preprint, Universiteit Utrecht and Universiteit Leiden,
2000.

\bibitem{Yoshida}
T.~Yoshida,
\emph{The relationship between face cuboids and elliptic curves},
preprint, 2024; arXiv:\allowbreak 2407.09825.

\end{thebibliography}

\end{document}
